\documentclass[11pt]{article}

\input{../shared/preamble.tex}

\usepackage{amsmath}
\usepackage{booktabs}
\usepackage{cleveref}

\title{A Tilt Is Worth What the Correlation Is Worth}
\author{Peter Cotton}
\date{September 21, 2026}

\newcommand{\E}{\mathbb{E}}
\newcommand{\R}{\mathbb{R}}
\newcommand{\one}{\mathbf{1}}

\newtheorem{proposition}{Proposition}
\newtheorem{remark}{Remark}

\begin{document}
\maketitle

\begin{abstract}
A long-only portfolio and a vector of latent abilities are the same object in
two charts: calibrate abilities so a race reproduces the portfolio, and the
portfolio is recovered by running the race. Re-running it under an estimated
correlation, damped by a parameter, tilts any benchmark without inverting
anything, so the whole dial is available where a covariance matrix is
singular. We measure what the dial is worth. It is worth nothing when the
correlation is noise and thirteen percent of variance when it is not: at a
tenth of an observation per asset a full tilt costs nineteen percent, and at
five observations per asset it saves thirteen, with the optimum walking
monotonically between. The effect replicates across three seeds and is
strongest where it is most significant, but it appears only on a market whose
covariance switches regime and is flat on a stationary one, which locates its
value in structure that no single covariance can hold. A tail-aware race,
which the construction admits and which the tilt is often motivated by, is
consistently worse than the Gaussian one on both variance and shortfall.
\end{abstract}

\section{The two charts}

Let $w$ be a long-only portfolio on the simplex interior. Treating its weights
as the winning probabilities of a race among assets, there is a vector of
latent abilities that reproduces them, and running the race at those
abilities returns the weights. Smaller ability means stronger, because the
winner is the minimum performance.

The two maps are inverse to each other, and a user can check it in one line.
Over random draws from three to sixty assets the round trip is accurate to
about $10^{-9}$, with one qualification stated in \cref{sec:scope}.

Nothing about that is a portfolio construction yet. It becomes one when the
race is re-run under a different dependence. Calibrate the abilities against a
reference correlation, blend that reference toward an estimated correlation by
a parameter $\phi$, and race again. At $\phi = 0$ the benchmark comes back. At
$\phi = 1$ the race runs under the full estimate. Nothing is inverted at any
point.

The construction therefore has the shape of a bridge, with a heuristic
benchmark at one end and a dependence-aware portfolio at the other, reached
without a matrix inverse. The question this paper answers is what the dial is
worth, and where.

\section{What the dial is worth}

The benchmark tilted is hierarchical risk parity, chosen because the tilt then
repairs a rule rather than competing with one. Measurements are paired under
common random numbers: the race at $\phi$ is compared against the same race at
$\phi = 0$ on the same fixed seed ensemble, so the simulation noise cancels
and what remains is the effect of the dial.

\begin{table}[ht]
\centering
\begin{tabular}{rrrl}
\toprule
$T/n$ & variance ratio & win rate & status \\
\midrule
0.10 & 1.19 & 27\% & tilting hurts \\
0.30 & 0.97 to 1.01 & 48 to 63\% & borderline \\
1.00 & 0.98 & 60 to 67\% & modest, replicates \\
2.00 & 0.90 to 0.92 & 82 to 84\% & replicated to a quarter of a standard error \\
5.00 & 0.87 & 88\% & strongest, one run \\
\bottomrule
\end{tabular}
\caption{Full tilt against no tilt, paired, forty assets, three seeds.}
\label{tab:dial}
\end{table}

The curve is monotone in coverage and that is the result. The dial behaves as
a trust parameter should: with four observations for forty assets the
estimated correlation is close to noise, tilting toward it costs nineteen
percent, and the correct setting is near zero. With five observations per
asset the same tilt saves thirteen percent of variance and four percent of
expected shortfall in roughly nine draws out of ten.

Expected shortfall follows the same shape rather than a different one. It is
neutral at parity, where the first run's 70 percent fell to 50 on
replication, and reaches 80 to 84 percent once there is real data.

\section{It needs something to express}

The same paired test on a stationary Gaussian market is flat: ratios between
0.99 and 1.01 with win rates from 47 to 53 percent at three tenths coverage.
The gains in \cref{tab:dial} come from a market whose covariance switches, a
calm structure plus a crash regime firing eight percent of the time with an
independent dependence structure over a permuted labelling.

That difference is the mechanism. The race is a non-linear functional of the
correlation, so it can express dependence that no single covariance matrix
holds. On a market that really is one Gaussian it has nothing to express, and
it does not pretend otherwise.

\section{The tail-aware race does not earn its keep}

The construction admits a Student-$t$ race, with fat marginals and genuine
tail dependence, and that is often the motivation for preferring a race to a
solve. It does not help.

On the crash market, scored on the same out-of-sample panel, the Gaussian
tilt beats the $t$ tilt at every sample ratio and on both objectives. At three
tenths coverage the Gaussian half-tilt realizes 0.979 of the benchmark's
variance against 1.016 for the $t$, and 0.995 of its shortfall against 1.009.

The market is not the reason. A multivariate $t$ market is elliptical, so
expected shortfall is proportional to standard deviation there and a tail
metric is the variance metric relabelled; the market used here is
non-elliptical by construction, with shortfall varying seventeen percent
relative to standard deviation across portfolios where the elliptical case
varies one percent.

\section{Scope}\label{sec:scope}

Three limits.

The round trip is not exact everywhere. The solver used for the inverse enters
a limit cycle on small, highly concentrated fields, and neither more
iterations, a tighter tolerance nor finer quadrature moves it. A damped
refinement clears the catastrophic cases, and what remains announces itself
rather than returning a wrong answer quietly. On four hundred concentrated
six-asset draws, fifteen warn and none is wrong in silence.

The strongest cell has one run behind it. Everything at two observations per
asset and below is replicated across seeds, and the replication was worth
running: it pulled all four shared cells down and eliminated one outright,
though every one of them passed a two-sigma check. The five-observation row
has not had that treatment.

The markets are simulated. The population covariance has to be known for
realized risk to be exact and for the elliptical check to be possible at all.

\section{Closing}

The dial is worth measuring because its answer is not constant. Tilting a
benchmark toward an estimated correlation is harmful when the estimate is
noise, valuable when it is not, and the crossover is at roughly a third of an
observation per asset on these markets. That is the behaviour a trust
parameter should have, and it is the same shape as the coupling parameter of
the Schur-complementary family \parencite{cotton2024schur}, reached by a
construction that never inverts a matrix.

Two results point away from the obvious extensions. The tilt is flat on a
stationary market, so it is not a general-purpose improvement but a way to act
on structure a covariance cannot hold. And the tail-aware race, which is the
natural thing to reach for next, is worse than the Gaussian one even on a tail
metric and even on a market with real asymmetric tail dependence.

\printbibliography

\end{document}
